\documentclass[a4paper]{article}

%Tutti gli usepackage vanno qui

\usepackage{geometry}
\usepackage[italian]{babel}
\usepackage[utf8]{inputenc}
\usepackage[T1]{fontenc}
\usepackage{tgschola}
%\usepackage{tgbonum}
\usepackage{tabularx}
\usepackage{longtable}
\usepackage{hyperref}
\usepackage{enumitem}
\usepackage[toc]{appendix}
\hypersetup{
	colorlinks=true,
	linkcolor=black,
	filecolor=magenta,      
	urlcolor=blue,
}
% Numerazione figure 
\let\counterwithout\relax
\let\counterwithin\relax
\usepackage{chngcntr}

\counterwithin{table}{subsection}
\counterwithin{figure}{subsection}

\usepackage[bottom]{footmisc}
\usepackage{fancyhdr}
\setcounter{secnumdepth}{4}
\usepackage{amsmath, amssymb}
\usepackage{array}
\usepackage{graphicx}

\usepackage{ifthen}

%\usepackage{float}
\usepackage{layouts}
\usepackage{url}
\usepackage{comment}
\usepackage{float}
\usepackage{eurosym}

\usepackage{lastpage}
\usepackage{layouts}
\usepackage{float}
\usepackage{eurosym}

%Comandi di impaginazione uguale per tutti i documenti
\pagestyle{fancy}
\lhead{\includegraphics[scale=0.04]{../../../../latex/images/logoTWM.png}}
%Titolo del documento
\rhead{\doctitle{}}
%\rfoot{\thepage}
\cfoot{Pagina \thepage\ di \pageref{LastPage}}
\setlength{\headheight}{35pt}
\setcounter{tocdepth}{5}
\setcounter{secnumdepth}{5}
\renewcommand{\footrulewidth}{0.4pt}

% multirow per tabelle
\usepackage{multirow}

% Permette tabelle su più pagine
%\usepackage{longtable}


% colore di sfondo per le celle
\usepackage[table]{xcolor}

%COMANDI TABELLE
\newcommand{\rowcolorhead}{\rowcolor[HTML]{9b240a}} %intestazione 
% check for missing commands
\newcommand{\headertitle}[1]{\textbf{\color{white}#1}} %titolo colonna
\definecolor{pari}{HTML}{FFDBCB}
\definecolor{dispari}{HTML}{F1F7FD}

% comandi glossario
\newcommand{\glo}{$_{G}$}
\newcommand{\glosp}{$_{G}$ }


%label custom
\makeatletter
\newcommand{\uclabel}[2]{%
	\protected@write \@auxout {}{\string \newlabel {#1}{{#2}{\thepage}{#2}{#1}{}} }%
	\hypertarget{#1}{#2}
}
\makeatother

%riportare pezzi di codice
\definecolor{codegray}{gray}{0.9}
\newcommand{\code}[1]{\colorbox{codegray}{\texttt{#1}}}



% Configurazione della pagina iniziale
\newcommand{\doctitle}{Norme di Progetto}
\newcommand{\docdate}{ } % lasciare vuoto (uno carattere di spazio) per i documenti che non sono verbali, così non viene scritta la data
\newcommand{\rev}{3.0.1}
\newcommand{\stato}{Non Approvato}
\newcommand{\uso}{Interno}
\newcommand{\approv}{--De Renzis Simone}
\newcommand{\red}{Greggio Nicolò\\& Tessari Andrea}
\newcommand{\ver}{Zuccolo Giada\\& Crivellari Alberto}
\newcommand{\dest}{Three Way Milkshake\\ Prof. Vardanega Tullio\\ Prof.
 Cardin Riccardo}
\newcommand{\describedoc}{Questo documento contiene tutte le norme di progetto\textsubscript{G}, definite inizialmente o aggiunte in seguito, viene quindi aggiornato per incrementi successivi.}



 % modifica questo file
\makeindex

\usepackage{ulem}
\usepackage{hyperref}


\begin{document}
	\thispagestyle{empty}
\begin{titlepage}
	\begin{center}
		
		\includegraphics[scale = 0.17]{../../../../latex/images/logoTWM.png}\\[0.7cm]
		

		\noindent\rule{\textwidth}{1pt} \\[0.4cm]
		\Huge \textbf{\doctitle} \\[0.1cm]
		\ifthenelse{\equal{\docdate}{ }}{ }{ \huge \textbf{\docdate} \\[0.1cm] }
		
		\noindent\rule{\textwidth}{1pt}\\[0.7cm]
		
		\large \textbf{Three Way Milkshake - Progetto "PORTACS"} \\[0.4cm] 
                \texttt{threewaymilkshake@gmail.com} \\[0.4cm]
                
		
        
        
        \large

        \begin{tabular}{r|l}
                        \textbf{Versione} & \rev{} \\
                        \textbf{Stato} & \stato{} \\
                        \textbf{Uso} & \uso{} \\                         
                        \textbf{Approvazione} & \approv{} \\                      
                        \textbf{Redazione} & \red{} \\ 
                        \textbf{Verifica} &  \ver{} \\                         
                        \textbf{Destinatari} & \parbox[t]{5cm}{ \dest{} }
                \end{tabular} 
                \\[0.3cm]
                \large \textbf{Descrizione} \\ \describedoc{} 
               

	\end{center}
\end{titlepage}
	\pagebreak

	% Registro delle modifiche
	\section*{Registro delle modifiche}

\newcommand{\changelogTable}[1]{
	
	
	\renewcommand{\arraystretch}{1.5}
	\rowcolors{2}{pari}{dispari}
	\begin{longtable}{ 
			>{\centering}p{0.07\textwidth} 
			>{}p{0.21\textwidth}
			>{\centering}p{0.17\textwidth}
			>{\centering}p{0.13\textwidth} 
			>{\centering}p{0.17\textwidth} 
			>{\centering}p{0.13\textwidth} }
		\rowcolorhead
		\headertitle{Vers.} &
		\centering \headertitle{Descrizione} &	
		\headertitle{Redazione} &
		\headertitle{Data red.} & 
		\headertitle{Verifica} &
		\headertitle{Data ver.}
		\endfirsthead	
		\endhead
		
		#1
		
	\end{longtable}
	\vspace{-2em}
	
}


\newcommand{\approvingTable}[1]{ 
	
	
	\renewcommand{\arraystretch}{1.5}
	\rowcolors{2}{pari}{dispari}
	\begin{longtable}{ 
			>{\centering}p{0.07\textwidth} 
			>{\centering}p{0.415\textwidth}
			>{\centering}p{0.13\textwidth}
			>{\centering}p{0.322\textwidth}  }
		\rowcolorhead
		\headertitle{Vers.} &
		\centering \headertitle{Descrizione} &	
		\headertitle{Data appr.} &
		\headertitle{Approvazione}
		\endfirsthead	
		\endhead
		
		#1
		
	\end{longtable}
	\vspace{-2em}
	
}
	\changelogTable{
	3.5.0 & Incremento appendice \S\ A e B & Crivellari Alberto & 2021-04-26 & -- & --
	3.4.0 & Rimossa appendice \S\ D & Crivellari Alberto & 2021-04-01 & -- & --
	\tabularnewline
	3.3.0 & Modifica sezioni \S\ 2 e \S\ 3 & Crivellari Alberto & 2021-04-01 & -- & --
	\tabularnewline
	3.2.0 & Modifica appendice \S\ B & Zuccolo Giada & 2021-03-31 & -- & --
\tabularnewline
	3.1.0 & Aggiornamento appendice \S\ C & Crivellari Alberto & 2021-03-20 & -- & --
}

\approvingTable{
    3.0.0 & Approvazione del documento & 2021-03-08 & Zuccolo Giada
}

\changelogTable{
    2.1.0 & Aggiornamento sezioni \S\ 4 e 5 & Crivellari Alberto & 2021-03-08 & Chiarello Sofia & 2021-03-08
}

\approvingTable{
    2.0.0 & Approvazione del documento & 2021-02-22 & De Renzis Simone
}

\changelogTable{
1.3.0 & Completamento stesura appendice \S\ D e \S\ B & Crivellari Alberto & 2021-02-12 & Greggio Nicolò & 2021-02-19
\tabularnewline
1.2.1 & Completamento stesura sezione \S\ 2 e \S\ 3 & Chiarello Sofia & 2021-02-11 & Tessari Andrea & 2021-02-13
\tabularnewline
1.2.0 & Stesura sezione \S\ 3 e appendice \S\ C & Chiarello Sofia & 2021-02-09 & Tessari Andrea & 2021-02-12
\tabularnewline
1.1.0 & Stesura sezione \S\ 2 & Chiarello Sofia & 2021-02-08 & Tessari Andrea & 2021-02-13
}

\approvingTable{
	1.0.0 & Approvazione del documento & 2021-01-10 & De Renzis Simone
}

\changelogTable{
0.5.1 & Aggiunta tabelle a sezione \S 4 e \S5 & Crivellari Alberto & 2021-01-08 & Greggio Nicolò & 2021-01-10
\tabularnewline
0.5.0 & Redazione sezione \S 5 & Crivellari Alberto & 2020-12-07 & Greggio Nicolò & 2021-01-10
\tabularnewline
0.4.0 & Redazione sezione \S 4 & Crivellari Alberto & 2020-12-06 & Greggio Nicolò & 2021-01-10
\tabularnewline
0.3.2 & Modifiche sezione \S 1 & Crivellari Alberto & 2020-12-30 & Greggio Nicolò & 2021-01-10
\tabularnewline
0.3.1 & Tabelle sezione \S 3 & Crivellari Alberto & 2020-12-29 & Greggio Nicolò & 2021-01-10
\tabularnewline
0.3.0 & Redazione sezione \S 3 & Tessari Andrea & 2020-12-28 & Greggio Nicolò & 2021-01-10
\tabularnewline
0.2.1 & Tabelle sezione \S 2 & Crivellari Alberto & 2020-12-20 &  Greggio Nicolò & 2021-01-09
\tabularnewline
0.2.0 & Redazione sezione \S 2 & Crivellari Alberto & 2020-12-19 & Greggio Nicolò & 2021-01-09
\tabularnewline
0.1.1 & Modifiche sezione \S 1  & Crivellari Alberto & 2020-12-18 & Greggio Nicolò & 2021-01-09
\tabularnewline
0.1.0 & Redazione sezione \S 1 & Crivellari Alberto & 2020-12-16 & Greggio Nicolò & 2021-01-09
\tabularnewline
0.0.1 & Strutturazione del documento & Tessari Andrea & 2020-12-15 & Greggio Nicolò & 2021-01-09
}


 % modifica questo file
	\pagebreak

	% indice
    {
        \hypersetup{linkcolor=black}
        \tableofcontents
    }
	\pagebreak

	% indice delle figure
	%\listoffigures
	%\pagebreak

	% indice delle tabelle
	%\listoftables
	%\pagebreak

	% contenuto del documento, ogni sezione in un file

	\section{Introduzione}
\subsection{Scopo del documento}
    Questo documento ha lo scopo di fissare e definire tutte le regole, convenzioni e buone pratiche utili a formare un way of working condiviso alla base da tutti i componenti del gruppo per assicurare una collaborazione efficiente ed efficace. Si discuteranno inoltre i vari strumenti che verranno adottati per facilitare lo sviluppo del progetto\textsubscript{G} e per promuovere un'organizzazione adeguata.
    Si ritiene inoltre che la definizione ed il mantenimento per incremento di un documento condiviso all'interno del gruppo di lavoro, che definisca e raccolga quanto descritto in maniera formale e centralizzata, possa favorire, in un contesto dove i membri possano variare, l'inserimento di nuovi componenti facilitandone l'ambientamento. Pur non essendo questo il contesto di lavoro attuale, è comunque una buona pratica da sperimentare e consolidare.

\subsection{Scopo del prodotto}
Il capitolato\textsubscript{G} C5 propone un progetto\textsubscript{G} in cui viene richiesto lo sviluppo di un software per il monitoraggio in tempo reale di unità che si muovono in uno spazio definito. All'interno di questo spazio, creato dall'utente per riprodurre le caratteristiche di un ambiente reale, le unità dovranno essere in grado di circolare in autonomia, o sotto il controllo dell'utente, per raggiungere dei punti di interesse posti nella mappa.  La circolazione è sottoposta a vincoli di viabilità e ad ostacoli propri della topologia dell'ambiente, deve evitare le collisioni con le altre unità e prevedere la gestione di situazioni critiche nel traffico.

\subsection{Termini, abbreviazioni ed altri documenti}
    Tutti i termini che necessitano di una spiegazione, per fornire un'adeguata comprensione, o perché possono causare ambiguità nel contesto, sono definiti nel glossario. A dispetto di tutti gli altri documenti presenti, il glossario è sotto la veste di una pagina web, più nello specifico in una wiki\textsubscript{G} di GitHub, accessibile a \href{https://github.com/Three-Way-Milkshake/docs/wiki/Glossario}{questo indirizzo}, con vocaboli suddivisi tra "Acronimi" e "Termini", disposti in ordine alfabetico al fine di facilitare la navigazione. All'interno dei documenti le voci di glossario saranno seguite da una G pedice mentre gli acronimi da una A (e.g.: voce di glossario\textsubscript{G} ; acronimo\textsubscript{A}).
    Inoltre quando si farà riferimento ad un altro documento o al documento stesso, il nome di questo sarà in maiuscoletto (e.g.: \textsc{esempio nome documento}).

\subsection{Riferimenti}
\label{ref}
    \subsubsection{Riferimenti Normativi}
        \begin{itemize}
            \item Standard ISO 12207:\\ \url{https://www.math.unipd.it/~tullio/IS-1/2009/Approfondimenti/ISO_12207-1995.pdf};
            \item Standard UML\textsubscript{A} 2.0:\\\url{https://www.omg.org/spec/UML/2.0/Superstructure/PDF};
            \item Diagrammi dei casi d'uso\textsubscript{G}:\\\url{https://www.math.unipd.it/~rcardin/swea/2021/Diagrammi\%20Use\%20Case_4x4.pdf};
            \item Diagrammi delle classi:\\\url{https://www.math.unipd.it/~rcardin/swea/2021/Diagrammi\%20delle\%20Classi_4x4.pdf};
            \item Diagrammi dei package:\\\url{https://www.math.unipd.it/~rcardin/swea/2021/Diagrammi\%20dei\%20Package_4x4.pdf};
            \item Diagrammi di attività\textsubscript{G}:\\\url{https://www.math.unipd.it/~rcardin/swea/2021/Diagrammi\%20di\%20Attivit\%c3\%a0_4x4.pdf};
            \item Diagrammi di sequenza:\\\url{https://www.math.unipd.it/~rcardin/swea/2021/Diagrammi\%20di\%20Sequenza_4x4.pdf};
            \item Standard ISO 8601:\\\url{https://www.iso.org/iso-8601-date-and-time-format.html}.
        \end{itemize}

    \subsubsection{Riferimenti Informativi}
        \begin{itemize}
        	\item \textsc{\href{https://github.com/Three-Way-Milkshake/docs/wiki/Glossario}{Glossario}}:\\per la definizione dei termini (pedice G) e degli acronimi (pedice A) evidenziati nel documento;
            \item \url{https://www.math.unipd.it/~tullio/IS-1/2020/Dispense/L03.pdf}
            \item \url{https://www.math.unipd.it/~tullio/IS-1/2020/Dispense/L06.pdf}
            \item \url{https://www.math.unipd.it/~tullio/IS-1/2020/Dispense/FC2.pdf}
            \item \url{https://www.math.unipd.it/~rcardin/swea/2021/SOLID\%20Principles\%20of\%20Object-Oriented\%20Design_4x4.pdf}
            \item \url{https://www.math.unipd.it/~tullio/IS-1/2020/Dispense/L09.pdf}
            \item \url{https://docs.github.com/en/free-pro-team@latest/actions}
            \item \url{https://www.latex-project.org/}
            \item \url{http://www.texstudio.org/}
            \item \url{https://www.xm1math.net/texmaker/}
            \item \url{https://www.atlassian.com/git/tutorials/comparing-workflows/gitflow-workflow}
            \item \url{https://github.com/about}
            \item \url{https://www.atlassian.com/software/jira}
            \item \url{https://www.atlassian.com/}
            \item \url{https://meet.google.com/}
            \item \url{https://slack.com/intl/en-it/about}
            \item \url{https://telegram.org/}
            \item \url{https://workspace.google.com/products/chat/}
        \end{itemize}

\pagebreak
 % andrea
	\pagebreak

	\section{Analisi del capitolato\textsubscript{G} scelto}
	\subsection{Capitolato C5 - Portacs}
\subsubsection{Informazioni generali}
	\begin{itemize}
	\item \textbf{Nome:} PORTACS\textsubscript{A};
	\item \textbf{Proponente:} Sanmarco Informatica;
	\item \textbf{Committente:}  Prof. Vardanega Tullio e Prof. Cardin Riccardo;
	\end{itemize}
\subsubsection{Descrizione del capitolato}
Il capitolato\textsubscript{G} C5 propone lo sviluppo di un software per il monitoraggio in tempo reale di varie unità e dei loro spostamenti in uno spazio creato dall'utente. \\
Ogni unità (rappresentabile da un robot cameriere, da un muletto o da un'auto a guida autonoma) dovrà raggiungere dei "Points Of Interest" (POI).
Ognuna di esse avrà: 
\begin{itemize}
	\item \textbf{Codice identificativo};
	\item \textbf{Velocità massima};
	\item \textbf{Posizione iniziale};
	\item \textbf{Lista ordinata dei POI\textsubscript{A} da attraversare}.
\end{itemize}
Lo spazio sarà definito da una mappa inserita dall'utente e avrà la struttura di una scacchiera con dei vincoli sulla viabilità (percorrenze in parallelo e/o sensi unici) e ostacoli in base alla topologia fornita.
Il software inoltre dovrà occuparsi di prevedere situazioni critiche, come il traffico ed eventuali collisioni esterne.

\subsubsection{Finalità del progetto}
Il progetto\textsubscript{G} PORTACS\textsubscript{A} si pone come obiettivo finale di dimostrare la fattibilità dello sviluppo di un software che permetta il coordinamento di unità che si muovono in uno spazio per raggiungere una lista ordinata di punti d'interesse. Questo implica che bisognerà sviluppare un progetto\textsubscript{G} negli ambiti del real time monitoring e analysis, e del predictivity e real time decision making. Inoltre, esso affronta le problematiche del mondo della logistica e dell'ottimizzazione delle performance di consegna.
Verranno quindi consegnati i seguenti documenti:
\begin{itemize}
	\item Codice sorgente di quanto realizzato;
	\item Docker file con la componente applicativa, rappresentante il motore di calcolo;
	\item Docker file con la componente applicativa, rappresentante il visualizzatore/monitor real time;
	\item Docker file, da istanziare N volte, rappresentante la singola unità;
	\item Docker file, da istanziare N volte, rappresentante il singolo pedone (facoltativo).
\end{itemize}
\subsubsection{Tecnologie interessate}
Le tecnologie interessate sono:
\begin{itemize}
	\item \textbf{Diagrammi UML};
	\item \textbf{\href{https://github.com/about}{Github}} o \textbf{\href{https://bitbucket.org/product/guides/getting-started/overview}{Bitbucket}};	
	\item \textbf{\href{https://www.docker.com/why-docker}{Docker}}.
\end{itemize}
\subsubsection{Aspetti positivi}
L'azienda mette a disposizione figure di diverso livello in modo tale da poter coprire nella maniera più appropriata le esigenze degli studenti.
Inoltre il codice prodotto sarà reso disponibile al pubblico con licenza libera su Github o BitBucket alla fine del progetto\textsubscript{G}.
\subsubsection{Criticità e fattori di rischio}
Per lo sviluppo di questo progetto\textsubscript{G}, si andrà a lavorare con il multithreading e sarà necessario prendere delle decisioni per il paradigma di programmazione che si sceglierà di usare. 
Tutto questo, sommato al fatto che si deve anche individuare un algoritmo che risolvi il problema del Path Finding, potrebbe non essere immediato e portare al riscontro di diverse difficoltà.
\subsubsection{Conclusioni}
La disponibilità dell'azienda assieme alla curiosità scaturita dopo il seminario offerto hanno favorito la scelta di questo capitolato\textsubscript{G}, già presa in considerazione visto l'argomento di interesse comune del gruppo riguardo gli argomenti di sviluppo proposti.
Inoltre questo progetto\textsubscript{G} aiuta a mettere in pratica le proprie conoscenze di ricerca operativa ed è infine molto utile acquisire competenze nel mondo Docker e nell'ambito del real time monitoring e analysis, e del predictivity e real time decision making. % andrea e un po' tutti
	\pagebreak

	\section{Analisi dei capitolato\textsubscript{G} rimanenti}
	\subsection{Capitolato C1 - BlockCOVID}


\subsubsection{Informazioni generali}

\begin{itemize}
	\item{\textbf{Nome:}} BlockCOVID
	\item{\textbf{Proponente:}} Imola Informatica
	\item{\textbf{Committente:}} Prof. Vardanega Tullio, Prof. Cardin Riccardo
\end{itemize}



\subsubsection{Descrizione del capitolato}

Il software BlockCOVID si pone l'obiettivo di fornire un'infrastruttura che consenta il tracciamento del personale e della pulizia delle postazioni di lavoro all'interno di un'azienda. Mira infatti a favorire l'attuazione e amministrazione delle misure necessarie per la tutela della salute dei lavoratori nel contesto della pandemia da COVID-19.



\subsubsection{Finalità del progetto}

Il prodotto finale prevede due modalità di tracciamento:
\begin{itemize}
	\item{\textbf{Tracciamento delle presenze:}} registrazione e monitoraggio in tempo reale delle presenze all'interno di un ambiente di lavoro (laboratorio informatico). Gli utenti devono avere la possibilità, in base alle disponibilità di postazioni, di effettuare prenotazioni tramite un applicativo per dispositivi mobili \textit{Android} e \textit{IOS}. Tramite lo stesso dispotivo segnaleranno la propria presenza una volta occupata la postazione ed eventualmente la sanificazione della stessa dopo l'uso. L'amministratore deve poter creare e gestire la struttura dell'ambiente, e visualizzarne in tempo reale lo stato di occupazione e sanificazione delle postazioni; 
	\item{\textbf{Tracciamento della pulizia delle postazioni:}} dedicato agli addetti alle pulizie, il prodotto deve consentire di visualizzare le postazioni non sanificate e segnalare le postazioni sanificate in seguito alla pulizia delle stesse.
\end{itemize}



\subsubsection{Tecnologie interessate}

\textit{Java} (versione 8 o superiori), \textit{Python} o \textit{Node.js}.
\begin{itemize}
	\item{\textbf{Java, Python, Node.js:}} linguaggi suggeriti dal proponente per lo sviluppo del server back-end che amministra l'infrastruttura di tracciamento;
	\item{\textbf{Protocolli asincroni:}} da utilizzare per gestire la comunicazione tra le componenti del sistema: i dispositivi mobili degli utenti, dell'amministratore e degli addetti alle pulizie con il server centrale;
	\item{\textbf{Blockchain:}} sistema che assicura la registrazione dei dati relativi alle sanificazioni conferendone valore legale;  
	\item{\textbf{Kubernetes IAAS, PAAS, Openshift, Rancher:}} piattaforme suggerite dal proponente per il rilascio delle componenti del server e la gestione della scalabilità orizzontale.
\end{itemize}

Inoltre il proponente specifica i seguenti requisiti minimi:
\begin{itemize}
	\item Fornire delle API che utilizzino tecnologie \textit{Rest} o \textit{gRPC} nell'ambito della comunicazione tra server e dispositivi;
	\item Studiare la tecnologia \textit{RFID} o eventuali alternative da utilizzare per certificare la presenza di persone nelle postazioni, analizzandone il consumo energetico;
	\item Effettuare test di integrazione e unità (copertura $\geq 80\%$) delle componenti applicative;
	\item Assicurare l'integrità del sistema tramite test end-to-end.
\end{itemize}
	
	
	
\subsubsection{Aspetti positivi}

Il progetto si pone in un contesto quanto mai attuale e di grande interesse generale. Trova infatti soluzione a problematiche diffuse e richieste da molte realtà, fornendo un caso di studio molto interessante per quanto riguarda i sistemi distribuiti e real-time. Le tecnologie impiegate e richieste sono valide, e il loro studio consentirebbe ai membri del gruppo di acquisire competenze facilmente spendibili nel mondo del lavoro.



\subsubsection{Criticità}

I requisiti del proponente sono dettagliati e appaiono abbastanza esigenti; le tecnologie da impiegare, per quanto interessanti, configurano un carico di lavoro consistente.

In generale, il capitolato ha riscosso un discreto interesse da parte del gruppo ma non si è configurato come prima scelta a favore di capitolati riguardanti temi considerati più interessanti, tecnologie più stimolanti e carichi di lavoro più sostenibili. 



\subsubsection{Conclusioni}

Nonostante l'elevato interesse e attenzione che questi sistemi hanno attirato, il gruppo ha deciso di procedere a favore di altri capitolati.
Per quanto attuale, il progetto si presuppone essere completato ed eventualmente utilizzato in un periodo in cui le problematiche legate alla pandemia dovrebbero, auspicabilmente, risultare fortemente mitigate dall'avvento delle campagne vaccinali. Non è comunque da escludere che, visto l'enorme impatto di questa emergenza, forme di tracciamento non debbano continuare ad essere applicate anche in condizioni di normalità.




 % simone
	\pagebreak

	\subsection{Capitolato C2 - Emporio Lambda}


\subsubsection{Informazioni generali}

\begin{itemize}
	\item{\textbf{Nome:}} Emporio Lambda
	\item{\textbf{Proponente:}} Red Babel
	\item{\textbf{Committente:}} Prof. Vardanega Tullio, Prof. Cardin Riccardo
\end{itemize}



\subsubsection{Descrizione del capitolato}

L'idea alla base di Emporio Lambda è quella di costruire una piattaforma di e-commerce che si basi interamente su tecnologie serverless.


\subsubsection{Finalità del progetto}

Il prodotto in questione dovrà fornire due insiemi di funzionalità principali: 
\begin{itemize}
	\item una orientata ai clienti (pagina principale, lista e descrizione prodotti, carrello degli acquisti, pagamento, gestione account);
	\item l'altra alle funzioni di back office (tutto ciò che può servire agli impiegati, ossia gestione di contabilità, inventario, ordini, giacenza, distribuzione, spedizioni...).
\end{itemize}

L'esecuzione del sistema avviene su architetture serverless che possono incorporare BaaS di terze parti, o che contengono codice proprietario eseguito su container effimeri (che hanno una durata limitata ad una singola invocazione) su piattaforme FaaS. SPA è l'approccio da adottare per la realizzazione dell'applicazione web. Tramite queste caratteristiche, il sistema nel suo complesso dovrebbe beneficiare di una riduzione in termini di costi di operazioni, complessità e tempi di consegna.


\subsubsection{Tecnologie interessate}

Il linguaggio principale che dev'essere adottato è \href{https://www.typescriptlang.org/}{Typescript}. Segue un elenco di altre tecnoglogie a supporto, obbligatorie o consigliate dal proponente:
\begin{itemize}
    \item \textbf{\href{https://aws.amazon.com/it/lambda/}{AWS Lambda}}: piattaforma di calcolo serverless basata su eventi, parte della suite di servizi web forniti da Amazon;
    \item \textbf{\href{https://aws.amazon.com/it/cloudformation/}{CloudFormation}}: strumento consigliato per il rilascio di infrastrutture AWS;
    \item \textbf{Serverless Framework}: framework Node.js da utilizzare per la parte back end;
    \item \textbf{\href{https://nextjs.org/}{Next.js}:} framework da usare per la parte front end;
    \item \textbf{\href{https://aws.amazon.com/it/cloudwatch/}{CloudWatch} o \href{https://www.datadoghq.com/}{Datadog}}: per implementare il sistema di monitoring;
    \item \textbf{\href{https://www.contentful.com/}{Contentful}}: CMS suggerito per l'implementazione di una parte opzionale.
\end{itemize}


\subsubsection{Aspetti positivi}

Emporio Lambda rappresenta sicuramente una realtà familiare a tutti, in quanto sono sempre più numerose e diffuse le piattaforme per gli acquisti online. Inoltre offre l'opportunità di lavorare con tecnologie nuove e all'avanguardia, adottando paradigmi diversi a quelli tradizionali e fornendo delle ottime opportunità di apprendimento.


\subsubsection{Criticità}

I requisiti imposti sembrano essere molto vincolanti, e le possibilità di scelte libere attuabili in fase di sviluppo appaiono limitate. Le piattaforme di e-commerce sono un dominio estremamente diffuso e conosciuto e, almeno dal lato dell'utente, non rappresentano una novità per i membri del team: altri capitolati sono stati giudicati più stimolanti.


\subsubsection{Conclusioni}

Questo capitolato non ha suscitato fin dall'inizio interesse nel gruppo. Non essendoci inoltre stato un seminario di approfondimento, non è stato possibile rivedere in chiave diversa le tematiche coinvolte, dunque l'opinione generale interna è rimasta invariata. Il gruppo si orienta verso un'altra scelta.




 % modifica questo file
	\pagebreak

	\subsection{Capitolato C3 - GDP - Gathering Detection Platform}


\subsubsection{Informazioni generali}

\begin{itemize}
	\item{\textbf{Nome:}} GDP - Gathering Detection Platform
	\item{\textbf{Proponente:}} SyncLab
	\item{\textbf{Committente:}} Prof. Vardanega Tullio, Prof. Cardin Riccardo
\end{itemize}



\subsubsection{Descrizione del capitolato}
In seguito al lockdown generato a causa del virus COVID-19, si sono venute a creare situazioni, lungo l'arco della vita quotidian,a in cui le persone entrano in contatto tra di loro, generando assembramento.\\
L’informatica può aiutare a governare meglio questa situazione.\\
Il software GDP (Gathering Detection Platform) è costituito da una piattaforma che rappresenta, mediante visualizzazione grafica, le zone potenzialmente a rischio di assembramento col fine di prevenirne di nuove, attraverso sensori e varie sorgenti di dati.\\


\subsubsection{Finalità del progetto}

Il software finale prevede l'acquisizione di informazione da sensoristica e da altre sorgenti:
\begin{itemize}
	\item{\textbf{Sensori:}}
	    \begin{itemize}
	        \item telecamere
	        \item dispositivi contapersone
	        \item etc.
	    \end{itemize}
	\item{\textbf{Sorgenti varie e eterogenee:}}
	    \begin{itemize}
	        \item flussi di prenotazioni Uber
	        \item orari dei mezzi di trasporto con capienze medie per corsia (autobus, metro, treno)
	        \item etc.
	\end{itemize}
\end{itemize}
Gli utilizzatori del software potranno vedere una rappresentazione delle zone a rischio (attuali o possibilmente in futuro), attraverso una heat-map.\\
In questo modo potranno accedere alla situazione globale dei vari flussi:
\begin{itemize}
    \item in tempo reale, con bassa latenza
    \item flussi previsti in futuro.
    \item flussi vecchi raccolti e storicizzati.
\end{itemize}


\subsubsection{Tecnologie interessate}

\begin{itemize}
	\item{\textbf{Java, \href{https://angular.io/}{Angular:}}} linguaggi suggeriti dal proponente per lo sviluppo del server back-end e della componente Web Application del sistema;
	\item{\textbf{\href{https://leafletjs.com/}{framework Leaflet}:}} framework\textsubscript{G} da utilizzare per la gestione delle mappe, ad esempio heatmap;
	\item{\textbf{protocolli asincroni:}} per la comunicazione tra le varie componenti;
	\item{\textbf{pattern Publisher/Subscriber e protocollo \href{https://mqtt.org/}{MQTT}:}}\newline (Message Queue Telemetry Transport), consigliato per essere open, di facile implementazione, di ampia diffusione in applicazioni IoT e M2M.
\end{itemize}

Inoltre il proponente specifica i seguenti requisito\textsubscript{G} minimi:
\begin{itemize}
	\item \textbf{responsive} deve sempre rispondere a una richiesta di servizio, anche in caso di guasto;
	\item \textbf{resilient} I servizi devono essere sempre ripristinabili a seguito di guasti;
	\item \textbf{elastic} I servizi devono scalare in base alla domanda;
	\item \textbf{message-driven} I servizi devono rispondere al mondo, senza controllarlo.
\end{itemize}
Un sistema che soddisfa questi requisito\textsubscript{G} è definito un Sistema Reattivo.


\subsubsection{Aspetti positivi}

Il progetto\textsubscript{G} si pone in un contesto molto tecnologico, di grande interesse per l'attuale mercato, molto sensibile e maturo su questa problematica.\\
Anche le tecnologie impiegate sono attuali e valide, il cui studio porterebbe all'acquisizione di competenze molto utili in ambito lavorativo.



\subsubsection{Criticità}

Nonostante il capitolato\textsubscript{G} abbia riscosso un buon interesse da parte del gruppo, le tecnologie da impiegare, in quanto alcune nuove e diverse da qualsiasi altro linguaggio imparato durante il percorso di studi, gravano molto sulla facilità di realizzazione del progetto\textsubscript{G}.\\
In particolare le criticità principali sono:
\begin{itemize}
    \item \textbf{Leaflet:} Tecnologia relativa alle heat-map e framework\textsubscript{G} Leaflet;
    \item \textbf{Linguaggio Angular:} Linguaggio di programmazione che viene eseguito su web browser;
    \item \textbf{Protocolli asincroni, Pattern Publisher/Subscriber e protocollo MQTT: } Varie tecnologie, alcune sconosciute, che servono alla realizzazione del progetto\textsubscript{G}.
\end{itemize}


\subsubsection{Conclusioni}

Nonostante l'elevata attenzione che questi sistemi hanno attirato, il gruppo non ha preso in considerazione questo capitolato\textsubscript{G} in quanto molti altri gruppi di lavoro avevano già segnalato un alto interesse per questo progetto\textsubscript{G}. \\
 % modifica questo file
	\pagebreak

	\subsection{Capitolato C4 - HD Viz, visualizzazione di dati con molte dimensioni}
\subsubsection{Informazioni generali}
	\begin{itemize}
	\item \textbf{Nome:} HD Viz, visualizzazione di dati con molte dimensioni;
	\item \textbf{Proponente:} Zucchetti;
	\item \textbf{Committente:}  Prof. Vardanega Tullio e Prof. Cardin Riccardo;
	\end{itemize}
\subsubsection{Descrizione del capitolato}
Il capitolato propone la realizzazione di un'applicazione chiamata "HD Viz", utile per l'interpretazione dell'analisi di dati a molte dimensioni, raccolti sfruttando tecniche di statistica, machine learning e intelligenza artificiale.
Richiede quindi lo sviluppo di un'applicazione per la visualizzazione di dati con molte dimensioni (ovvero i dati da visualizzare dovranno poter avere almeno fino a 15 dimensioni e deve però essere possibile visualizzare anche dati con meno dimensioni) a supporto della fase esplorativa di questi dati da interpretare, forniti al sistema di visualizzazione sia con query che da file in formato CSV.
La visualizzazione di dati con il software "HD Viz" deve essere presentata almeno con i grafici "Scatter plot Matrix", "Force Field", "Heat Map" e "Proiezione Lineare Multi Asse". In particolare dovrà ordinare i punti nel grafico "Heat Map" per evidenziare i "cluster" presenti nei dati.
Inoltre sono proposti altri requisiti valutati come opzionali.
\subsubsection{Finalità del progetto}
L'obiettivo è sviluppare un'applicazione per la visualizzazione di dati, di molte dimensioni, raccolti con la finalità di conseguire poi con un'analisi dei dati.
Il software "HD Viz" dovra presentare almeno le visualizzazioni con:
\begin{itemize}
\item \textbf{Scatter plot Matrix} (fino ad un massimo di 5 dimensioni);
\item \textbf{Force Field};
\item \textbf{Heat Map} (obbligatoriamente dovrà organizzare i punti per evidenziare i "cluster" presenti nei dati);
\item \textbf{Proiezione Lineare Multi Asse};
\end{itemize}
\subsubsection{Tecnologie interessate}
L'azienda indica lo sviluppo dell'applicazione "HD Viz" con:
\begin{itemize}
\item \textbf{HTML/CSS/JavaScript} utilizzando la \textit{libreria D3.js};
\item \textbf{Java con server Tomcat} o \textbf{JavaScript con server Node.js} per la parte server di supporto alla presentazione nel browser e alle query ad un database SQL o NoSQL.
\end{itemize}
\subsubsection{Aspetti positivi}
Questo capitolato propone lo sviluppo un progetto che si inserisce in un ambito relativamente nuovo, cioè la raccolta, la gestione, la memorizzazione e la visualizzazione di grandi quantità di dati, i Big Data, al giorno d'oggi utilizzabili in tutti i vari campi.
Lo sviluppo di questo progetto avrebbe permesso l'approfondimento di questo nuovo ma più complesso ambito, utilizzando delle tecnologie non totalmente nuove alla gran parte dei membri del gruppo.
\subsubsection{Criticità e fattori di rischio}
Sebbene le tecnologie interessate per lo sviluppo dell'applicazione non siamo totalmente nuove alla gran parte dei membri del gruppo, applicarle per il soddisfacimento delle richieste obbligatorie e dei requisiti opzionali, potrebbe portare ad una situazione svantaggiosa in quanto, sia le richieste obbligatorie sia i requisiti opzionali, siano più complessi da sviluppare e nel complesso anche numerose.
\subsubsection{Conclusioni}
Malgrado l'ambito abbastanza stimolante, il capitolato ha riscosso un misurato interesse da parte del gruppo.
Per questo motivo si è scelto di selezionare altri capitolati che riscontravano maggiore curiosità nel gruppo e che si avvalgono di tecnologie nuove. % modifica questo file
	\pagebreak

	\subsection{Capitolato C6 - RGP: Realtime Gaming Platform}
\subsubsection{Informazioni generali}
	\begin{itemize}
	\item \textbf{Nome:} RGP: Realtime Gaming Platform;
	\item \textbf{Proponente:} Zero12 s.r.l.;
	\item \textbf{Committente:}  Prof. Tullio Vardanega e Prof. Riccardo Cardin;
	\end{itemize}
\subsubsection{Descrizione del capitolato}
    Il capitolato proposto prevede la realizzazione di un videogioco a scorrimento verticale, fruibile da dispositivi mobile, con la possibilità di giocare in multiplayer real time.
    La modalità di gioco, per quanto riguarda il gameplay, sarà simile ad Aero Fighters\footnote{Esempio Aero Fighters su youtube: \url{https://youtu.be/5Phj-735p30?t=28}}, mentre la grafica potrà essere scelta liberamente dal gruppo o fornita dal team di Zero12.

\subsubsection{Finalità del progetto}
    Sono previste le modalità di gioco single e multiplayer, ma
    il cuore del progetto sono le sfide tra più giocatori: il gioco è ad eliminazione, l'ultimo rimasto vince.
    Durante la partita deve essere possibile vedere, in tempo reale, i movimenti del rivale e per garantire l'equità bisognerà prevedere la sincronizzazione di eventuali nemici e powerup.
    La modalità single player invece consiste in una serie infinita di livelli a difficoltà crescente. Il gioco termina quando l'utente ha concluso le vite oppure se non ha raccolto power-up sufficienti a mantenere il proprio oggetto attivo.
    Il progetto è finalizzato allo sviluppo di un'applicazione mobile superando dei vincoli quali:
    \begin{itemize}
    	\item ricerca delle tecnologie AWS per capire quale si può adattare meglio ad un gioco con requisiti di realtime, raccogliendo le motivazioni a supporto della scelta;
    	\item realizzazione della componente server-side;
    	\item implementazione del gioco per piattaforma mobile.
    \end{itemize}
\subsubsection{Tecnologie interessate}
    Le tecnologie interessate sono:
    \begin{itemize}
    	\item \textbf{AWS: }per la realizzazione dei servizi per la gestione dei giochi;
    	\item \textbf{NodeJs: }per lo sviluppo di codice server side;
    	\item \textbf{Swift o Kotlin: }per l'implementazione delle applicazioni mobile, rispettivamente per iOS e Android;
    	\item \textbf{Git: }per il versionamento del codice sviluppato.
    \end{itemize}
\subsubsection{Aspetti positivi}
    Il capitolato permette di entrare nel mondo dello sviluppo mobile nativo e della gestione di comunicazioni real time fra più dispositivi. Inoltre consente anche di familiarizzare con le tecnologie serverless anche tramite i servizi di Amazon AWS, per i quali l'azienda proponente prevede corsi di formazione.
\subsubsection{Criticità e fattori di rischio}
    L'ambiente AWS è molto vasto ed è critica la scelta iniziale del servizio sul quale basarsi non avendo alcuna esperienza pregressa: una scelta errata può portare ad una realizzazione non ottimale, rischiando così di non soddisfare pienamente i requisiti real time della parte multiplayer del videogioco. Inoltre lo sviluppo di applicazioni mobile native è già un compito arduo di per sé, per cui aggiungere comunicazioni real time e sincronizzazioni tra dispositivi diversi aumenta ulteriormente il grado di difficoltà e le possibilità d'errore.
\subsubsection{Conclusioni}
    I corsi di formazione offerti dall'azienda potrebbero essere interessanti, ma valutata l'elevata difficoltà e lo scarso interesse da parte del gruppo, la scelta non si pone qui. % modifica questo file
	\pagebreak

	\subsection{Capitolato C7 - Soluzioni di sincronizzazione Desktop}
\subsubsection{Informazioni generali}
	\begin{itemize}
	\item \textbf{Nome:} \textit{Soluzioni di sincronizzazione Desktop;}
	\item \textbf{Proponente:} \textit{Zextras;}
	\item \textbf{Committente:}  \textit{Prof. Tullio Vardanega e Prof. Riccardo Cardin;}
	\end{itemize}
\subsubsection{Descrizione del capitolato}
Il capitolato richiede di sviluppare un algoritmo di sincronizzazione Desktop solido ed efficiente in grado di garantire il salvataggio in cloud del lavoro, in modo da poter raggiungere in qualsiasi momento e da qualsiasi dispositivo il proprio lavoro. Inoltre deve essere sviluppata un'interfaccia multipiattaforma per l'uso dell'algoritmo nei più importanti sistemi operativi esistenti (MacOs, Windows, Linux) senza richiedere all'utente l'installazione di ulteriori prodotti per il funzionamento.
\subsubsection{Finalità del progetto}
L'obiettivo è quello di creare questo algoritmo di sincronizzazione che funzioni in più piattaforme in grado di garantire il salvataggio in cloud del lavoro e contemporaneamente la sincronizzazione dei cambiamenti presenti in cloud. Inoltre viene richiesto l'utilizzo dell’algoritmo sviluppato per richiedere e fornire i cambiamenti ai contenuti in sincronizzazione verso il prodotto Zextras Drive.
L'algoritmo deve avere le seguenti principali funzionalità:
\begin{itemize}
\item Configurazione ed autenticazione dell’utente;
\item Gestione di cosa sincronizzare e di cosa ignorare nelle cartelle cloud;
\item Gestione di cosa sincronizzare e di cosa ignorare nelle cartelle locali;
\item Sincronizzazione costante dei cambiamenti, siano essi locali o remoti;
\item Possibilità di modifica delle preferenze a posteriori;
\item Sistema di notifica utente dei cambiamenti.
\end{itemize}
\subsubsection{Tecnologie interessate}
L'azienda esprime la necessità di sviluppare l'algoritmo per i più importanti sistemi operativi esistenti (MacOs, Windwos, Linux) e consiglia l'utilizzo:
\begin{itemize}
\item \textbf{Qt Framework:} strumento basato sul linguaggio C++ orientato ad oggetti, da utilizzare per creare l'interfaccia poiché fortemente supportato e documentato;
\item \textbf{Python:} strumento consigliato per lo sviluppo della Business Logic, linguaggio ad alto livello adatto allo sviluppo di applicazioni distribuite.
\end{itemize}
\subsubsection{Aspetti positivi}
Il progetto si pone in un contesto molto utilizzato sia da utenti privati che da aziende. Inoltre le tecnologie consigliate fanno già parte delle conoscenze di gran parte dei membri del gruppo e ci sono molti esempi, anch'essi ben conosciuti, a cui ispirarsi per la creazione dell'algoritmo (Dropbox, Google Drive, ect).
\subsubsection{Criticità e fattori di rischio}
Il confronto con tecnologie già esistenti di questo tipo può essere svantaggioso poiché si potrebbe creare un prodotto molto più inefficiente. Inoltre le richieste sono numerose e abbastanza complesse.
\subsubsection{Conclusioni}
Si è deciso di puntare su altri capitolati poiché questo non ha suscitato grande interesse per quanto riguarda la tematica che affronta e il settore interessato. % modifica questo file
	\pagebreak


\end{document}
